\begin{table}
\centering
\caption{Scoring functions $\psi(e_s, e_r, e_o)$ of KG embedding methods. $e_s, e_r, e_o \in \Re^d$ for TransE, DistMult, ConvE, InteractE, and ComDensE, $e_s, e_r, e_o \in \mathrm{C}^d$ for ComplEx and RotatE. $\circ$ denote Hadamard-product, * denote convolution, $\bullet$ denote circular convolution, $\Psi(\mathcal{P}_k)$ denote feature permutation, $f$ and $g$ denote activation functions, $[\bar{e_s};\bar{e_r}]_{2d}$ represent 2d reshaping, $[e_s;e_r]_{1d}$ represent vector concatenation, $\Omega$ denote shared dense layer and $\Omega_r$ indicate relation-aware operation.}
\label{table1}
\begin{tabular}{c|c}
    \textbf{Model} & \textbf{Scoring function $\psi(e_s, e_r, e_o)$}\\
    \hline
    TransE & $||e_s + e_r - e_o||_p$\\
    RotatE & $-||e_s \circ e_r - e_o||^2$\\ 
    DistMult & $<e_s, e_r, e_o>$\\
    ComplEx & $Re(<e_s, e_r, e_o>)$\\
    ConvE & $f(vec(f([\bar{e_s};\bar{e_r}]_{2d} * w))W)e_o$\\
    InteractE & $g(vec(f(\Psi(\mathcal{P}_k) \bullet w))W)e_o$\\
    \hline
    ComDensE & $f([f(\Omega_r([e_s;e_r]_{1d})); f(\Omega([e_s;e_r]_{1d}))]_{1d}^TW)e_o$
\end{tabular}
\end{table}

\begin{table*}[t]
\centering
\caption{Statistics of benchmark datasets. PageRank of datasets is cited from~\cite{dettmers2018convolutional}.}
\label{table2}
\begin{tabular}{c|c|c|c|c|c|c}
 \textbf{Dataset} &
 \textbf{\#Entities} &
 \textbf{\#Relations} &
 \textbf{\#Training triples} &
 \textbf{\#Validation triples} &
 \textbf{\#Test triples} &
 \textbf{PageRank} \\
 \hline
 FB15k-237 & 14,541 & 237 & 272,115 & 17,535 & 20,466 & 0.733 \\
 WN18RR & 40,943 & 11 & 86,835 & 3,034 & 3,134 & 0.104
 \\
\end{tabular}
%}
\end{table*}

\section{Method}
This section describes the details of ComDensE architecture which is depicted in Fig.~\ref{fig2}. Let $e_s, e_o \in \Re^{d_e}$ denote the embedding vectors of subject and object entities, and $e_r \in \Re^{d_r}$ denote the embedding vector of relations. In order to compute the scoring function of ComDensE, we begin by concatenating the subject and relation vectors as follows:
\begin{equation}
[e_s;e_r]_{1d} \in \Re^{d}  
\label{eq1}
\end{equation}
where $d := d_e + d_r$ and $[a;b]_{1d}$ denote vector concatenation of vectors $a$ and $b$. The concatenated embedding is the input to all the subsequent layers.

\subsection{Common Feature Extraction Layer}
The common feature extraction is performed by a dense layer as well. This layer is commonly applied to all types of input embeddings. The width of a dense layer corresponds to the number of filters of the layer where each filter contains kernel size equal to that of input embeddings. The width is the key design parameter, in which we explain in more details in the experiment section.  

In the shared dense layer, we first apply an affine function $\Omega(\cdot)$ to input embeddings given by 
\begin{equation}
\Omega(x) = {W_h} x + b_h
\label{eq2}
\end{equation}
where $W_h \in \Re^{n d_h \times d}$ and $b_h \in \Re^{nd_h}$. The width of dense layer is given by $nd_h$, a multiple of $d_h$ where $n$ is the hyperparameter to be determined. The output of common feature extraction is obtained by applying the nonlinear activation $f(\cdot)$ to $\Omega([e_s;e_r]_{1d})$.

\subsection{Relation-aware Feature Extraction Layer}
In order to extract relation-specific features from concatenated embeddings (\ref{eq1}), we consider the relation-aware encoding function. The encoding function is denoted by $\Omega_r$ for relation $r$. Note that, if any triple contains the relation given by $r$, $\Omega_r(\cdot)$ is applied to input embeddings irrespective of subject or object entities. 
In ComDensE, we use a dense layer for the relation-aware feature extraction. The encoding function $\Omega_r$ is an affine function given by
\begin{equation}
\Omega_r(x) = {W_r}x + b_r
\label{eq3}
\end{equation}
where $W_r \in \Re^{d_z\times d } $, $b_r \in \Re^{d_z}$ and $d_z$ is the output length of $\Omega_r$. We first apply $\Omega_r$ to input embeddings (\ref{eq1}), and then apply the nonlinear activation function $f(\cdot)$.

\subsection{Projection to Embedding Space}
After obtaining the latent vectors from relation-aware and common feature extraction layers, the vectors are concatenated. The concatenated vector is projected to the embedding space by projection matrix $W \in \Re^{(d_z+nd) \times d_e}$, and then nonlinear activation $f(\cdot)$ is applied. Specifically, let us define $h_{sr} \in \Re^{d_e}$ as 
\begin{equation}
h_{sr} := f([f(\Omega_r([e_s;e_r]_{1d})); f(\Omega([e_s;e_r]_{1d}))]_{1d}^TW)
\label{eq4}
\end{equation}
The link prediction score $\psi(e_s, e_r, e_o)$ is defined as the inner product $h_{sr}^Te_o$. See Table \ref{table1} for the comparison of scoring function $\psi$. 

\subsection{Loss Function}
We compute the scores for all the triple, and calculate the loss using binary cross entropy function. We use 1:N training strategy introduced by~\cite{dettmers2018convolutional}. Let $N$ denote the number of all entities in KG. The loss $\mathcal{L}$ is given by
\begin{equation}
\mathcal{L} = -{1 \over N} \sum_{i=1}^N [(y_i) \log (p_i) + (1-y_i) \log (1-p_i)]
\label{eq5}
\end{equation}
where $p_i = \sigma(h_{sr}^Te_o^{(i)})$, $e_o^{(i)}$ is the $i$-th object entity, $y_i \in \left\{ 0,1 \right\}$ is the label, and $\sigma$ denote the sigmoid function.

\subsection{Properties of ComDensE model}

\noindent\textbf{Relation-specific inductive bias is important.} Typically in KGs, the number of relations is less than the number of entities, e.g., the ratio is 0.03\% for WN18RR, 1.6\% in FB15k-237, 0.03\% in YAGO3-10. This implies that triples in KGs having different relations tend to share multiple entities. Thus, relational inductive bias is needed so as to capture discriminative features of each relation in the presence of a large number of entities shared among different relations. ComDensE leverages both relation-aware and common feature extraction from embeddings to get the best of both approaches.

\noindent\textbf{Full-sized filters for concatenated embeddings are effective.} It was noted earlier that, it is important to mix as many components from different types of embeddings, i.e., entity and relation types when filtering the concatenated embeddings. Specifically in InteractE~\cite{vashishth2020interacte}, the authors propose the concept of  {\em heterogeneous interaction}. The heterogeneous interaction is defined as the degree to which a local (convolutional) filter can spatially cover different types, i.e., entity and relation types, of embedding components in the 2-D shaped concatenated embeddings. InteractE showed that the link prediction performance can be improved by increasing the heterogeneous interaction. Clearly, ComDensE maximizes such heterogeneous interaction by construction, because its dense layers use full-sized filters which entirely cover the 2-D area of the concatenated (and reshaped) embedding vectors of entities and relations.

\noindent\textbf{Computational overhead is manageable.} One of the merits of convolutional layers is parameter efficiency.  Indeed, the parameter count of dense layers is higher than the convolutional layers with smaller kernel sizes. We however argue that this incurs manageable computational overhead for KG embeddings. The size of KG embeddings are on the order of hundreds, which is much smaller than high-resolution image features. Furthermore, we later increase the width to capture abundant relational information for improving link prediction accuracy. In practice, prediction accuracy is often prioritized over parameter efficiency. Experiments also show that DensE incurs reasonable parameter counts compared to previous methods.